% =====================================================================
%  附录 A 李代数与泛包络代数
%  覆盖：PDF p.389–p.395（原书印刷 p.379–p.385）
%  A.1–A.5；公式一律以 pages/p389–p395.png 为唯一依据重排
% =====================================================================

\chapter{李代数与泛包络代数}

\section{Lie 代数}

\begin{keypoints}
    \item \textbf{核心思想：}李代数是带双线性括号的向量空间，括号满足反称性与 Jacobi 恒等式。
    \item \textbf{数学工具：}括号 $[\cdot,\cdot]:\mathfrak g\otimes\mathfrak g\to\mathfrak g$；表示是满足 $\rho([x,y])=\rho(x)\rho(y)-\rho(y)\rho(x)$ 的线性映射 $\rho:\mathfrak g\to\End(V)$。
    \item \textbf{关键结论：}李代数是李群在单位元处切空间的线性化；本书用李代数语言，因为圈群与仿射 Kac--Moody 代数的表示论主要在线性层面展开。
    \item \textbf{记号提醒：}全书采用 $[x,y]=-[y,x]$ 与循环和形式的 Jacobi 恒等式，后文所有结构常数均以此为准。
\end{keypoints}

% -------------------------- 印刷 p.379（PDF p.389） --------------------------

\textbf{李代数}是带双线性运算（即李括号）
\[
    [\cdot,\cdot]:\mathfrak g\otimes\mathfrak g\to\mathfrak g
\]
的向量空间 $\mathfrak g$，它满足两个附加条件：
\begin{itemize}
    \item $[x,y]=-[y,x]$；
    \item $[x,[y,z]]+[y,[z,x]]+[z,[x,y]]=0$。
\end{itemize}
后一条件称为 \textbf{Jacobi 恒等式}。

% [Tutor Note: 易错（p.379）。本书的符号约定是 $[x,y]=-[y,x]$、Jacobi 恒等式写作三项循环和为零；后文（A.3）的 Serre 关系、$[h_i,e_j]=a_{ij}e_j$、$[e_i,f_j]=\delta_{i,j}h_i$ 以及第 1、2 章的结构常数全部以此约定为准。若按相反约定 $[x,y]_{\text{新}}=-[x,y]_{\text{旧}}$ 阅读，所有 $a_{ij}$ 都要变号。]

李代数与李群密切相关。李群单位元处的切空间自然具有李代数的结构。因此，思考李代数是把来自李群论的问题线性化的一种方式。

% [Tutor Note: 含糊点拨（p.379）。“李群单位元处的切空间自然具有李代数结构”一句省略了构造：切空间 $T_1G$ 上由左不变向量场的交换子给出括号；本书只用这一对应来引入记号，不展开李群理论。对圈群 $G(\!(t)\!)$ 而言，对应的“李群”是形式群，因此本书自始至终在李代数与顶点代数层面工作。]

\textbf{李代数 $\mathfrak g$ 的表示}，或等价地，$\mathfrak g$-模，是一对 $(V,\rho)$，其中 $V$ 是向量空间，$\rho:\mathfrak g\to\End(V)$ 是满足
\[
    \rho([x,y])=\rho(x)\rho(y)-\rho(y)\rho(x)
\]
的映射。我们常常滥用记号，把表示称为 $V$ 而不写出同态 $\rho$；但 $\rho$ 始终隐含其中。

% [Tutor Note: 易错（p.379）。表示的定义要求 $\rho$ 是李代数同态，即 $\rho([x,y])=[\rho(x),\rho(y)]$，右端的括号是结合代数 $\End(V)$ 的交换子；这就是上式右端出现两项乘积之差的原因。第 2 章起“$\gghat$-模”的记号同此约定。]

\section{泛包络代数}

\begin{keypoints}
    \item \textbf{核心思想：}泛包络代数 $U(\mathfrak g)$ 是把李括号实现为交换子的“最自由”结合代数，其泛性质给出李代数同态到结合代数的唯一提升。
    \item \textbf{数学工具：}张量代数 $T(\mathfrak g)$、由 $g\otimes h-h\otimes g-[g,h]$ 生成的双边理想、PBW 滤过与分次伴随代数 $\operatorname{gr}$。
    \item \textbf{关键结论：}$\operatorname{gr}U(\mathfrak g)\cong\Sym(\mathfrak g)$（Poincaré--Birkhoff--Witt，定理 A.2.2）；作为向量空间 $U(\mathfrak g)\cong\Sym(\mathfrak g)$，但同构非典范。
    \item \textbf{记号提醒：}$\operatorname{gr}$ 指分次伴随代数，$\operatorname{Symb}$ 是符号映射；$U(\mathfrak g)_{\le i}$ 是 PBW 滤过的第 $i$ 项。
\end{keypoints}

% -------------------------- 印刷 p.379–380（PDF p.389–p.390） --------------------------

研究李代数及其表示的一个标准工具是泛包络代数。它是一个结合代数，可由任意李代数构造，并具有如下性质：李代数上的任意模都可视为该结合代数上的模。

由于我们希望在结合代数中把 $\mathfrak g$ 的元素相乘，考察 $\mathfrak g$ 的张量代数是有意义的：
\[
    T(\mathfrak g)=\C\oplus\mathfrak g\oplus\mathfrak g^{\otimes 2}\oplus\mathfrak g^{\otimes 3}\oplus\cdots .
\]

若给定 $\mathfrak g$ 的表示 $(V,\rho)$，我们显然希望相应的 $T(\mathfrak g)$-模使得 $T(\mathfrak g)$ 的作用在第二个直和项上限制为 $(V,\rho)$。又因为它应当是结合代数上的模，我们还希望
\[
    \rho(g\otimes h)=\rho(g)\rho(h).
\]
然而，李代数模还满足一个附加关系
\[
    \rho([g,h])=\rho(g)\rho(h)-\rho(h)\rho(g).
\]
因此 $T(\mathfrak g)$ 中的元素 $[g,h]$ 与 $g\otimes h-h\otimes g$ 应当以相同方式作用。这表明我们应在张量代数中把 $[g,h]$ 与 $g\otimes h-h\otimes g$ 等同起来，而这引出下述定义。

% [Tutor Note: 原书笔误（p.380）。原文 “we clearly want want the corresponding $T(\mathfrak g)$-module” 多写了一个 “want”；译文按语义只译一次，此处登记。]

% [Tutor Note: 背景（p.380）。泛性质等价于说 $U$ 是遗忘函子 $\Lie:\text{结合代数}\to\text{李代数}$ 的左伴随：$\Hom_{\mathrm{Lie}}(\mathfrak g,\Lie(A))\cong\Hom_{\mathrm{Ass}}(U(\mathfrak g),A)$。这就是“泛包络”中“泛”字的准确含义，也解释了为何 $U(\mathfrak g)$-模范畴与 $\mathfrak g$-模范畴等价。]

\textbf{泛包络代数} $U(\mathfrak g)$ 是张量代数 $T(\mathfrak g)$ 对由形如
\[
    g\otimes h-h\otimes g-[g,h]
\]
的元素生成的双边理想的商。注意有自然映射 $\mathfrak g\to U(\mathfrak g)$。

它之所以称为泛的，是因为具有下述泛性质：若 $A$ 是任意结合代数（自然地视为李代数），$f:\mathfrak g\to A$ 是李代数同态，则它可以唯一地补全为下述交换图：
\[
    \begin{array}{ccc}
         & U(\mathfrak g) & \\[-4pt]
         & \uparrow & \searrow \\[-4pt]
        \mathfrak g & \longrightarrow & A
    \end{array}
\]

由张量代数上的分次，泛包络代数上有一个自然 \textbf{滤过}（有时称为 PBW 滤过）。滤过的第 $i$ 项记为 $U(\mathfrak g)_{\le i}$。

当 $\mathfrak g$ 交换时，泛包络代数的结构很容易算出。此时理想的生成元就是 $g\otimes h-h\otimes g$（因为括号为零）。这恰好是由 $\mathfrak g$ 生成的对称代数的定义。于是得到
\[
    U(\mathfrak g)\cong\Sym(\mathfrak g).
\]
当 $\mathfrak g$ 非交换时，情形更复杂。

给定任意滤过代数
\[
    0=F_{\le-1}\subset F_{\le0}\subset F_{\le1}\subset\cdots\subset F,
\]
我们定义其 \textbf{分次伴随代数}为
\[
    \operatorname{gr}F=\bigoplus_{i\ge0}F_{\le i}/F_{\le i-1}.
\]
它自然是分次代数；作为向量空间有 $\operatorname{gr}F\cong F$，但这一同构并非典范。取商的分次伴随代数相当容易：
\[
    \operatorname{gr}(F/I)=\operatorname{gr}(F)/\operatorname{Symb}(I),
\]
其中 $\operatorname{Symb}$ 是 \textbf{符号映射}。设 $x$ 是滤过代数 $F$ 的元素。存在整数 $i$ 使得 $x\in F_{\le i}$ 而 $x\notin F_{\le i-1}$。$x$ 的符号定义为 $x$ 在商 $F_{\le i}/F_{\le i-1}$ 中的像。换言之，符号映射取出 $x$ 的首项。\transnote{“并非典范”是指这个向量空间同构依赖于在每个 $F_{\le i}$ 中选取 $F_{\le i-1}$ 的补空间（或选取代表元），不同选法给出不同同构。对 $U(\mathfrak g)$ 而言，$\operatorname{gr}U(\mathfrak g)\cong\Sym(\mathfrak g)$ 是分次代数同构（PBW 定理），但它只给出 $U(\mathfrak g)\cong\Sym(\mathfrak g)$ 的向量空间同构——若为代数同构，$U(\mathfrak g)$ 就交换了。此外，$\operatorname{Symb}(I)$ 是“$I$ 中元素的首项所生成的理想”，不是 $I$ 的像。[来源：原书 p.381。]}

% [Tutor Note: 易错（p.381）。$\operatorname{gr}F\cong F$ 只是向量空间同构，且非典范（原文 “but non-canonically”）；作商公式 $\operatorname{gr}(F/I)=\operatorname{gr}(F)/\operatorname{Symb}(I)$ 要求 $I$ 是滤过理想，且 $\operatorname{Symb}(I)$ 按分次取。符号映射取的是“首项”，不是整个元素。]

于是对泛包络代数有
\[
    \operatorname{gr}U(\mathfrak g)=\operatorname{gr}T(\mathfrak g)/\operatorname{Symb}(I).
\]
由于张量代数已分次，$\operatorname{gr}T(\mathfrak g)\cong T(\mathfrak g)$。

\begin{lemma}
    理想 $\operatorname{Symb}(I)$ 由理想 $I$ 的生成元的符号生成。
\end{lemma}

% [Tutor Note: 含糊点拨（p.381，引理 A.2.1）。作者只写下这一引理，随后说“这一事实的证明是 PBW 定理大多数证明的核心”而未给证明。要点是：$I$ 由 $g\otimes h-h\otimes g-[g,h]$ 生成，其符号为 $g\otimes h-h\otimes g$；要证这些符号生成整个分次理想，关键是 Jacobi 恒等式保证用 $\operatorname{ad}$ 反复作用生成理想时不会出现无法由生成元首项拼出的新首项。这正是 PBW 的实质，标准证明见 [Dixmier 1977]、[Kac 1990]。]

这一事实的证明是 Poincaré--Birkhoff--Witt 定理大多数证明的核心。它本质上依赖于括号满足 Jacobi 恒等式这一事实。如果我们用任意的双线性形式来定义理想 $I$（例如用由 $a\otimes b-b\otimes a-B(a,b)$ 生成的理想），那么 $\operatorname{Symb}(I)$ 中就会出现不能由生成元的符号得到的元素。生成元的符号容易算出，恰好是 $g\otimes h-h\otimes g$。相应的商是 $\Sym(\mathfrak g)$。于是我们证明了，对任意李代数 $\mathfrak g$ 都有
\[
    \operatorname{gr}U(\mathfrak g)\cong\Sym(\mathfrak g).
\]
换言之，作为向量空间，$U(\mathfrak g)$ 与 $\Sym(\mathfrak g)$ 同构，但这一同构并非典范。这常以下述形式引用。

\begin{theorem}[Poincaré--Birkhoff--Witt]
    设 $J^a$，$a=1,\dots,\dim\mathfrak g$，是 $\mathfrak g$ 的有序基（作为向量空间），则泛包络代数有一组基，由形如 $J^{a_1}\cdots J^{a_n}$（其中 $a_1\le\cdots\le a_n$）的元素构成。
\end{theorem}

% [Tutor Note: 易错（p.381，定理 A.2.2）。$J^a$ 是 $\mathfrak g$ 的“有序基”，单项式 $J^{a_1}\cdots J^{a_n}$ 的指标满足 $a_1\le\cdots\le a_n$，是弱递增（允许相等），不是严格递增；定理断言的是向量空间的基，不同次序的同一组因子在 $U(\mathfrak g)$ 中只相差低阶项。]

这种形式的单项式常称为\textbf{按字典序排列}的。\transnote{定理 A.2.2 中的 $a_1\leq\cdots\leq a_n$ 是弱递增（允许相等），这样每个单项式只出现一次；原书称这样的单项式为“按字典序排列”，是指把有序基 $\{J^a\}$ 固定次序后写成不降序列而言。注意 $U(\mathfrak g)$ 中不同次序的乘积并不相等，只相差低阶项，因此 PBW 基与 $\Sym(\mathfrak g)$ 的单项式基只是（非典范的）同构对应。[来源：原书 p.381 定理 A.2.2。]}研究泛包络代数的主要原因之一是：$U(\mathfrak g)$-模范畴与 $\mathfrak g$-模范畴等价（这不过是说，$\mathfrak g$ 的任意表示可赋予唯一的 $U(\mathfrak g)$-模结构）。

\section{单李代数}

\begin{keypoints}
    \item \textbf{核心思想：}有限维复单李代数由正定 Cartan 矩阵 $(a_{ij})$ 分类，共四个无限族 $A_n,B_n,C_n,D_n$ 与五个例外型 $E_6,E_7,E_8,F_4,G_2$。
    \item \textbf{数学工具：}Cartan 分解 $\mathfrak g=\mathfrak n_-\oplus\mathfrak h\oplus\mathfrak n_+$、单根 $\alpha_i$ 与单余根 $h_i$、Serre 关系。
    \item \textbf{关键结论：}权格 $P$、根格 $Q$、余权格 $\check P$、余根格 $\check Q$ 与外尔群 $W$ 由根数据决定；这些正是 §1.1.5 中构造 Langlands 对偶群的原料。
    \item \textbf{记号提醒：}$\alpha_j(h_i)=a_{ij}$ 定义单余根；$\check P\subset\mathfrak h$、$\check Q\subset\mathfrak h$ 而 $P,Q\subset\mathfrak h^*$，切勿混同。
\end{keypoints}

% -------------------------- 印刷 p.382（PDF p.392） --------------------------

本书考虑有限维复单李代数。这里“单”指李代数没有非平凡理想。这些李代数已被完全分类。它们与正定 \textbf{Cartan 矩阵} $A=(a_{ij})_{i,j=1,\dots,\ell}$ 一一对应。\transnote{“正定”排除的是双曲型等广义 Cartan 矩阵：正定对应半单李代数，而“单”还要求矩阵不可分解（等价于 Dynkin 图连通）。本书的约定是 $a_{ij}=\alpha_j(h_i)=\langle\alpha_j,\alpha_i^\vee\rangle=2(\alpha_i,\alpha_j)/(\alpha_i,\alpha_i)$，故 $a_{ij}$ 与 $a_{ji}$ 分别记录 $\alpha_j$ 在 $h_i$ 上的值与两端根长之比；$a_{ij}a_{ji}=4\cos^2(\pi/m_{ij})\in\{0,1,2,3\}$，其中 $m_{ij}$ 是 $s_is_j$ 的阶，只取 $2,3,4,6$。Dynkin 图就是这些整数的图形化：$i\neq j$ 时连 $|a_{ij}a_{ji}|$ 条边，$|a_{ij}|$ 与 $|a_{ji}|$ 孰大决定哪端是长根。[来源：原书 p.382（原文即引 [Kac 1990]）；Kac 1990, Ch.4（广义 Cartan 矩阵的分类）。]}这些矩阵是元素为整数的可对称化矩阵，满足以下条件：$a_{ii}=2$，且对 $i\neq j$ 有 $a_{ij}\le0$，并且 $a_{ij}=0$ 当且仅当 $a_{ji}\neq0$（详见 [Kac 1990]）。这样的矩阵，从而单李代数，归入四个无限族与六个例外型。四个族称为 $A_n$、$B_n$、$C_n$、$D_n$；例外李代数称为 $E_6$、$E_7$、$E_8$、$F_4$ 与 $G_2$。\transnote{此处原书有两处印刷错误，译文照录、在此更正：条件应为 $a_{ij}=0$ 当且仅当 $a_{ji}=0$（零元位置对称），原书印作 $a_{ji}\neq0$，与 $a_{ij}\leq0$ 合看自相矛盾；例外型应为“五个”而非“六个”，原文同段紧接着只列出 $E_6,E_7,E_8,F_4,G_2$，四个无限族加五个例外型正是 $A$--$G$ 型分类的全部。[来源：原书 p.382 同段；Kac 1990, 引言（广义 Cartan 矩阵要求 $a_{ij}=0\Rightarrow a_{ji}=0$）与 Ch.4（分类）。]}

% [Tutor Note: 原书笔误（p.382）。原文说分类为“四个无限族与六个例外型”，但紧接着列出的例外型只有五个（$E_6,E_7,E_8,F_4,G_2$）；标准分类是四个无限族加五个例外型。译文按原书保留“六个”，已登记报告。]

% [Tutor Note: 原书笔误（p.382）。条件应为 “$a_{ij}=0$ 当且仅当 $a_{ji}=0$”（Cartan 矩阵的零元位置对称），原文印成 “$a_{ij}=0$ 当且仅当 $a_{ji}\neq0$”，这样写自相矛盾。译文按原书保留，已登记报告。]

单李代数 $\mathfrak g$ 有 Cartan 分解
\[
    \mathfrak g=\mathfrak n_-\oplus\mathfrak h\oplus\mathfrak n_+.
\]
这里 $\mathfrak h$ 是 Cartan 子代数，即 $\mathfrak g$ 的极大交换李子代数，由半单元素（即其伴随作用在 $\mathfrak g$ 上半单的元素）组成。其维数为 $\ell$，即 $\mathfrak g$ 对应的 Cartan 矩阵的秩；我们也称 $\ell$ 为 $\mathfrak g$ 的\textbf{秩}。李子代数 $\mathfrak n_+$（相应地，$\mathfrak n_-$）称为上（相应地，下）幂零子代数。在 $\mathfrak h$ 的伴随作用下，幂零子代数分解为特征空间：
\[
    \mathfrak n_\pm=\bigoplus_{\alpha\in\Delta_+}\mathfrak n_{\pm\alpha},
\]
其中
\[
    \mathfrak n_{\pm\alpha}=\{x\in\mathfrak n_\pm\mid [h,x]=\pm\alpha(h)x\}.
\]
这里 $\Delta_+\subset\mathfrak h^*$，$\Delta_-=-\Delta_+$ 分别是 $\mathfrak g$ 的正根与负根 $\pm\alpha$ 的集合（对之 $\mathfrak n_{\pm\alpha}\neq0$）。事实上这些子空间都是一维的；我们分别取 $\mathfrak n_\alpha$ 与 $\mathfrak n_{-\alpha}$ 的生成元 $e_\alpha$ 与 $f_\alpha$。

% [Tutor Note: 易错（p.382）。下标要成对读：$\mathfrak n_+$ 对应正根 $\Delta_+$，$\mathfrak n_{\pm\alpha}$ 中的 $\pm$ 与根 $\pm\alpha$ 同步；$e_\alpha\in\mathfrak n_\alpha$ 而 $f_\alpha\in\mathfrak n_{-\alpha}$。若把 $\mathfrak n_+$ 与 $\mathfrak n_-$ 弄反，Serre 关系与 $[e_i,f_j]=\delta_{i,j}h_i$ 的符号会整体出错。]

已知 $\Delta_+$ 含有一族\textbf{单根} $\alpha_i$，$i=1,\dots,\ell$，其中 $\ell=\dim\mathfrak h$ 是 $\mathfrak g$ 的秩。$\Delta_+$ 的其余元素都是单根以非负整数系数的线性组合。单根构成 $\mathfrak h^*$ 的一组基。设 $h_i$，$i=1,\dots,\ell$，是由公式
\[
    \alpha_j(h_i)=a_{ij}
\]
唯一确定的 $\mathfrak h$ 的基。这些就是\textbf{单余根}。

我们取一维空间 $\mathfrak n_{\alpha_i}$（相应地，$\mathfrak n_{-\alpha_i}$）的非零生成元 $e_i=e_{\alpha_i}$（相应地，$f_i=f_{\alpha_i}$），$i=1,\dots,\ell$。则 $\mathfrak n_+$（相应地，$\mathfrak n_-$）由 $e_1,\dots,e_\ell$（相应地，$f_1,\dots,f_\ell$）生成，并满足下述 \textbf{Serre 关系}：
\[
    (\operatorname{ad} e_i)^{-a_{ij}+1}\cdot e_j=0,\qquad
    (\operatorname{ad} f_i)^{-a_{ij}+1}\cdot f_j=0,\qquad
    i\neq j.
\]

% [Tutor Note: 易错（p.382）。Serre 关系的指数是 $(\operatorname{ad} e_i)^{1-a_{ij}}$：因为 $a_{ij}\le0$，指数恰为 $1-a_{ij}\ge1$，即“对 $e_j$ 连续作用 $1-a_{ij}$ 次得零”。有些教材用对称化后的 $|a_{ij}|$ 或 $a_{ij}a_{ji}$ 表述，指数写法会差一个 1，务必以本书原文为准。]

% -------------------------- 印刷 p.383（PDF p.393） --------------------------

这里 $a_{ij}$ 是 $\mathfrak g$ 的 Cartan 矩阵的元素。

李代数 $\mathfrak g$ 由 $h_i$、$e_i$、$f_i$（$i=1,\dots,\ell$）生成，满足 Serre 关系以及
\[
    [h_i,e_j]=a_{ij}e_j,\qquad [h_i,f_j]=-a_{ij}f_j,\qquad [e_i,f_j]=\delta_{i,j}h_i.
\]
为满足最后一组关系，我们可能需要用非零数重新标度 $e_i$（或 $f_i$）。\transnote{原书此处印作“由 $h_i,h_i,f_i$ 生成”，第二个 $h_i$ 显系 $e_i$ 之误（译文按 $e_i$ 翻译）。下一句说可能需要重新标度 $e_i$（或 $f_i$），原因是：$h_i$ 已被 $\alpha_j(h_i)=a_{ij}$ 唯一固定，而在具体实现中 $[e_i,f_i]$ 可能只是 $h_i$ 的非零倍数，须把 $e_i$（或 $f_i$）乘非零数把系数归一到 $1$；这只改变生成元的标度，不改变由 Cartan 矩阵给出的抽象李代数。[来源：原书 p.383；Kac 1990, Ch.1（由 $h_i,e_i,f_i$ 与 Serre 关系定义 $\mathfrak g(A)$）。]}

% [Tutor Note: 原书笔误（p.383）。原文作 “generated by $h_i,h_i,f_i$”，第二个 $h_i$ 显系 $e_i$ 之误；译文已按 $e_i$ 翻译（公式未动），此处登记，并在正文对应句末出正式译注。]

例如，李代数 $A_n$ 就是所谓 $\mathfrak{sl}_{n+1}$。它是迹为零的 $(n+1)\times(n+1)$ 矩阵构成的向量空间，括号为 $[A,B]=AB-BA$。这是特殊线性群 $SL_{n+1}(\C)$（即行列式为 $1$ 的 $(n+1)\times(n+1)$ 矩阵构成的群）的李代数。为看出这一点，我们把 $SL_{n+1}(\C)$ 的元素写成 $I_{n+1}+\epsilon M$，其中 $\epsilon^2=0$，再看行列式条件对 $M$ 施加什么限制。这恰好就是迹为零的条件。

Cartan 子代数由对角矩阵组成。李代数 $\mathfrak n_+$ 与 $\mathfrak n_-$ 分别是上、下幂零子代数，其对角元为零。生成元 $e_i$、$h_i$、$f_i$ 分别是矩阵 $E_{i,i+1}$、$E_{ii}-E_{i+1,i+1}$ 与 $E_{i+1,i}$。

$\mathfrak g$ 上有四个格：权格 $P\subset\mathfrak h^*$ 由满足 $\varphi(h_i)\in\Z$ 的线性泛函 $\varphi:\mathfrak h\to\C$ 组成；根格 $Q\subset P$ 由单根 $\alpha_i$（$i=1,\dots,\ell$）在 $\Z$ 上张成；余权格 $\check P\subset\mathfrak h$ 由使单根取整数值的元素组成；余根格 $\check Q$ 由单余根 $h_i$（$i=1,\dots,\ell$）张成。\transnote{这四个格正是根数据的具体化：$P$、$\check P$ 分别对应特征标格 $X^*(T)$ 与余特征标格 $X_*(T)$（$T$ 取单连通群的极大环面），$Q$、$\check Q$ 分别由单根、单余根张成；再选定单根集合 $\Delta_s$ 即得 p.5--p.6 的带基根数据（based root datum）。要点是 $P,Q$ 在 $\mathfrak h^*$、$\check P,\check Q$ 在 $\mathfrak h$，配对 $\langle\cdot,\cdot\rangle:\mathfrak h^*\times\mathfrak h\to\C$ 是二者唯一的桥；Langlands 对偶把 $X^*$ 与 $X_*$ 互换，于是 $P$ 与 $\check P$、根与余根互换（$SL_n$ 与 $PGL_n$ 互换即最简例）。[来源：原书 p.5--p.6、p.383；并见本版 §1.1.5 译注。]}

格 $P$ 由满足 $\omega_i(h_j)=\delta_{i,j}$ 的元素 $\omega_i\in\mathfrak h^*$ 张成。余根格由满足公式 $\alpha_i(\check\omega_i)=\delta_{i,j}$ 的元素 $\check\omega_i\in\mathfrak h$ 张成。\transnote{本句原书有一处概念错误和一处指标错误：按上一段，余根格已定义为 $\check Q=\operatorname{span}_{\mathbb Z}\{h_i\}$，而满足 $\alpha_j(\check\omega_i)=\delta_{i,j}$ 的元素 $\check\omega_i$ 是基本余权，故主语应为“余权格”$\check P$，公式左端应为 $\alpha_j(\check\omega_i)$（原文两处指标都写作 $i$，与右端 $\delta_{i,j}$ 不闭合）。译文按原书保留“余根格”与 $\alpha_i(\check\omega_i)$ 的字面，读者按上述理解即可。[来源：原书 p.383（同页上一段的四格定义）；并见本版 §1.1.5（根数据中 $P$ 与 $\check P$ 的对称地位）。]}

% [Tutor Note: 原书笔误（p.383）。余根格生成元的公式应为 $\alpha_j(\check\omega_i)=\delta_{i,j}$（与权格的 $\omega_i(h_j)=\delta_{i,j}$ 对偶），原文左端两处指标都写成 $i$，右端却保留 $\delta_{i,j}$ 中的 $j$，指标不闭合。另：原书此处主语 “The coroot lattice” 亦误——上一段已把余根格定义为 $\operatorname{span}_{\mathbb Z}\{h_i\}$，而满足 $\alpha_j(\check\omega_i)=\delta_{i,j}$ 的元素是基本余权、张成余权格 $\check P$，故本句主语应为余权格。译文按原书保留，已登记报告。]

% [Tutor Note: 易错（p.383）。四个格分居两侧：$P,Q\subset\mathfrak h^*$（线性泛函），$\check P,\check Q\subset\mathfrak h$（其上有单根取整数值）；配对 $\langle\lambda,\check\alpha_i\rangle$ 是 $\mathfrak h^*\times\mathfrak h\to\C$ 的自然配对。把 $\check P$ 与 $P$ 对调是对偶群计算中最常见的错误来源（对照 §1.1.5）。]

我们注意，$P$ 可等同于与 $\mathfrak g$ 具相同李代数的连通单连通李群 $G$（例如 $SL_n$）的 Cartan 子群 $H_{\mathrm{sc}}$ 的特征标群 $H_{\mathrm{sc}}\to\C^\times$；而 $\check P$ 等同于余特征标群 $\C^\times\to H_{\mathrm{adj}}$，其中 $H_{\mathrm{adj}}$ 是与 $G$ 相伴的伴随型群（中心平凡，例如 $PGL_n$）的 Cartan 子群。

% [Tutor Note: 背景（p.383）。这一段正是 Langlands 对偶群的“字母表”：$P=\Hom(H_{\mathrm{sc}},\C^\times)$、$\check P=\Hom(\C^\times,H_{\mathrm{adj}})$，故交换 $X^*$ 与 $X_*$ 就是把 $P$ 与 $\check P$、单根与余根互换（§1.1.5 的 ${}^L G^\circ$ 由此而来）。$SL_n$ 与 $PGL_n$ 互为对偶，正是 $P$ 与 $\check P$ 互换的最简例子。]

设 $s_i$，$i=1,\dots,\ell$，是 $\mathfrak h^*$ 上由公式
\[
    s_i(\lambda)=\lambda-\langle\lambda,\check\alpha_i\rangle\alpha_i
\]
定义的线性算子（反射）。这些算子生成与 $\mathfrak g$ 相伴的外尔群 $W$ 在 $\mathfrak h^*$ 上的作用。\transnote{这里 $\check\alpha_i$ 是单余根的另一种写法，即 $\check\alpha_i=h_i=\alpha_i^\vee$，并非基本余权 $\check\omega_i$。用它可检验：$s_i(\alpha_i)=-\alpha_i$，且 $s_i(\alpha_j)=\alpha_j-a_{ij}\alpha_i$（$j\neq i$），因为 $a_{ij}=\alpha_j(h_i)$；这说明 $\{s_i\}$ 给出外尔群在权空间 $\mathfrak h^*$ 上的反射作用。[来源：原书 p.383；Kac 1990, Ch.6（外尔群与根系统）。]}

\section{中心扩张}

\begin{keypoints}
    \item \textbf{核心思想：}李代数的中心扩张是正合列 $0\to\mathfrak a\to\gghat\to\mathfrak g\to0$，其中 $\mathfrak a$ 落在中心里。
    \item \textbf{数学工具：}取向量空间分裂 $\gghat\cong\C\mathbf 1\oplus\mathfrak g$ 后，扩张由二元上循环 $c:\mathfrak g\otimes\mathfrak g\to\C$ 描述，满足反称性与 Jacobi 型条件。
    \item \textbf{关键结论：}一维中心扩张的同构类由 $H^2(\mathfrak g,\C)$ 分类，即二元上循环模上边缘。
    \item \textbf{记号提醒：}括号记为 $[x,y]_{\mathrm{new}}=[x,y]_{\mathrm{old}}+c(x,y)\mathbf 1$；$\mathbf 1$ 是中心生成元，$\iota_1,\iota_2$ 是分裂映射。
\end{keypoints}

% -------------------------- 印刷 p.383–384（PDF p.393–p.394） --------------------------

设 $\mathfrak a$ 是阿贝尔李代数。所谓李代数 $\mathfrak g$ 被 $\mathfrak a$ 的\textbf{中心扩张}，是指一个李代数 $\gghat$，它处于李代数的正合列
\begin{equation}
    0\to\mathfrak a\to\gghat\to\mathfrak g\to 0
    \tag{A.4-1}
\end{equation}
之中。这里 $\mathfrak a$ 是 $\gghat$ 的李子代数，且在 $\gghat$ 中中心，即对一切 $a\in\mathfrak a$、$x\in\mathfrak g$ 有 $[a,x]=0$。\transnote{严格说，“$\mathfrak a$ 在 $\gghat$ 中中心”要求 $[\mathfrak a,\gghat]=0$；原书这里把 $x$ 写成 $\mathfrak g$ 的元素，是因为 $\mathfrak a$ 阿贝尔且扩张的括号在 $\mathfrak a$ 上为零，检验 $x\in\mathfrak g$ 即可（$x\in\mathfrak a$ 时自动成立）。另外，下文的分裂 $\gghat\cong\C\mathbf 1\oplus\mathfrak g$ 只是向量空间同构：换一个分裂会改变上循环 $c$，但不改变扩张的同构类。[来源：原书 p.383--384；并见本节后文对同构类的分类。]}

% [Tutor Note: 易错（p.383，式 (A.4-1)）。中心扩张按定义要求 $\mathfrak a$ 落在 $\gghat$ 的中心：$[\mathfrak a,\gghat]=0$。正合列中的第一个箭头是包含，第二个是商映射；下文总取 $\mathfrak a=\C\mathbf 1$ 一维，此时 $\gghat$ 作为向量空间同构于 $\C\mathbf 1\oplus\mathfrak g$，但这一分裂不是典范的。]

设 $\mathfrak a\cong\C\mathbf 1$ 是一维的，生成元为 $\mathbf 1$。取 $\gghat$ 作为向量空间的分裂，$\gghat\cong\C\mathbf 1\oplus\mathfrak g$。则 $\gghat$ 上的李括号诱导出 $\C\mathbf 1\oplus\mathfrak g$ 上的李括号，使 $\C\mathbf 1$ 中心，且
\[
    [x,y]_{\mathrm{new}}=[x,y]_{\mathrm{old}}+c(x,y)\mathbf 1,
\]
其中 $x,y\in\mathfrak g$，$c:\mathfrak g\otimes\mathfrak g\to\C$ 是线性映射。上述公式定义 $\gghat$ 上的李括号，当且仅当 $c$ 满足：
\begin{itemize}
    \item $c(x,y)=-c(y,x)$；
    \item $c(x,[y,z])+c(y,[z,x])+c(z,[x,y])=0$。
\end{itemize}
这样的表达式称为 $\mathfrak g$ 上的\textbf{二元上循环}。

% [Tutor Note: 易错（p.384）。上循环的两个条件与式子 $[x,y]_{\mathrm{new}}=[x,y]_{\mathrm{old}}+c(x,y)\mathbf 1$ 的符号约定牢牢绑定：第一个条件来自反称性，第二个正是 Jacobi 恒等式；若把 $c$ 的符号改号（即写成 $[x,y]_{\mathrm{old}}-c(x,y)\mathbf 1$），下文的上边缘公式 $c_2=c_1+f([x,y])$ 也要相应变号。]

现在设给定形如 (A.4-1) 的两个扩张 $\gghat_1$ 与 $\gghat_2$，其中 $\mathfrak a=\C\mathbf 1$，并给定它们之间的一个保持各自子空间 $\C\mathbf 1$ 的李代数同构 $\phi$。取这些扩张的分裂 $\iota_1:\C\mathbf 1\oplus\mathfrak g\to\gghat_1$ 与 $\iota_2:\C\mathbf 1\oplus\mathfrak g\to\gghat_2$，得到二元上循环 $c_1$ 与 $c_2$。定义线性映射 $f:\mathfrak g\to\C$ 为 $\iota_1$、$\phi$ 与由 $\iota_2$ 诱导的投影 $\gghat_2\to\C\mathbf 1$ 的复合。于是容易看出
\[
    c_2(x,y)=c_1(x,y)+f([x,y]).
\]
因此，$\mathfrak g$ 的一维中心扩张的同构类集合（实际上是一个向量空间）同构于 $\mathfrak g$ 上二元上循环构成的向量空间对下述子空间的商：该子空间由满足
\[
    c(x,y)=f([x,y]),\qquad \forall x,y\in\mathfrak g,
\]
（对某个 $f:\mathfrak g\to\C$）的二元上循环 $c$ 组成。这样的上循环称为\textbf{上边缘}。该商空间是 $\mathfrak g$ 的第二李代数上同调群，记作 $H^2(\mathfrak g,\C)$。\transnote{对半单（特别地单）李代数，Whitehead 引理给出 $H^2(\mathfrak g,\C)=0$，所以本节的一维中心扩张对 $\mathfrak g$ 本身只给出平凡扩张；中心扩张真正的用武之地是 §A.5 的圈代数 $L\mathfrak g=\mathfrak g\otimes\C[t,t^{-1}]$ 及其形式版本，那里 $H^2$ 是一维的，并由不变内积 $\kappa$ 参数化。这也是全书从第 2 章起反复使用 $\kappa$ 的原因。[来源：原书 p.384；并见 §A.5 与第 2 章。]}

% [Tutor Note: 背景（p.384）。$H^2(\mathfrak g,\C)$ 分类一维中心扩张，这是李代数上同调的入门例子；本书真正用到的是 $\mathfrak g(\!(t)\!)$ 的中心扩张——它由 $H^2(\mathfrak g(\!(t)\!),\C)$ 分类，对单李代数是一维的，并由不变内积 $\kappa$ 参数化（见 A.5 与第 2 章）。临界水平 $\kappa_c$ 正是在这一维参数线上的特殊点。]

\section{仿射 Kac--Moody 代数}

\begin{keypoints}
    \item \textbf{核心思想：}把 $[\mathfrak g,\mathfrak g]$ 与交换结合代数 $A$ 张量，得到取值于 $\mathfrak g$ 的“函数”构成的李代数；取 $A=\C[t,t^{-1}]$ 得多项式圈代数，取 $A=\C((t))$ 得形式圈代数。
    \item \textbf{数学工具：}$[\mathfrak g\otimes a,\mathfrak h\otimes b]=[\mathfrak g,\mathfrak h]\otimes ab$；记号 $\mathfrak g\otimes_\C A$、$\C[t,t^{-1}]$、$\C((t))$。
    \item \textbf{关键结论：}中心扩张由 $H^2(L\mathfrak g,\C)$ 分类，对单李代数是一维的，由不变内积 $\kappa$ 参数化；形式版本给出仿射 Kac--Moody 代数 $\gghat_\kappa$。
    \item \textbf{记号提醒：}$L\mathfrak g$ 是多项式圈代数，$\gghat^{\mathrm{pol}}_\kappa$ 是它的中心扩张，$\gghat_\kappa$ 是形式圈代数 $\mathfrak g(\!(t)\!)$ 的中心扩张；本书主要用后者。
\end{keypoints}

% -------------------------- 印刷 p.384–385（PDF p.394–p.395） --------------------------

若给定复李代数 $\mathfrak g$ 与交换结合代数 $A$，则 $\mathfrak g\otimes_\C A$ 在括号
\[
    [g\otimes a,h\otimes b]=[g,h]\otimes ab
\]
下是李代数。

$A$ 究竟是什么？通常可以把交换结合代数看成某个流形 $M$ 上的函数。注意
\[
    \mathfrak g\otimes_\C\operatorname{Fun}(M)=\operatorname{Map}(M,\mathfrak g)
\]
正是从 $M$ 到 $\mathfrak g$ 的映射构成的李代数。它关于逐点李括号自然成为李代数（与上述公式给出的李括号一致）。

% [Tutor Note: 含糊点拨（p.384）。$\mathfrak g\otimes_\C\operatorname{Fun}(M)=\operatorname{Map}(M,\mathfrak g)$ 的等同要求把 $\operatorname{Fun}(M)$ 取为某个函数代数（多项式函数或代数函数），而不是任意交换结合代数；本书只取 $A=\C[t,t^{-1}]$ 或 $\C((t))$，因此不需要一般流形的函数论。]

让我们取流形为圆周，即考虑从圆周到李代数的映射。此时须决定处理哪一类映射。使用光滑映射需要一些分析，而我们希望停留在代数范围内——因此使用代数（多项式）映射。考虑嵌入复平面的单位圆，坐标为 $t$，于是在圆上 $t=\ee^{\ii\theta}$，$\theta$ 为角度。该圆上的多项式函数等同于 $t$ 的洛朗多项式。它们构成向量空间，记作 $\C[t,t^{-1}]$。用这个函数空间得到\textbf{多项式圈代数}
\[
    L\mathfrak g=\mathfrak g\otimes\C[t,t^{-1}],
\]
括号如上。

% [Tutor Note: 易错（p.385）。$t=\ee^{\ii\theta}$ 中的 $\theta$ 是圆周上的角度；$\C[t,t^{-1}]$ 是圆上的（洛朗）多项式函数代数，而 $\C((t))$ 是形式洛朗级数代数。前者是几何对象（圆上的环路），后者是形式对象（去心形式圆盘），两者只在代数层面相近；本书从第 2 章起统一使用 $\C((t))$，这正是“形式圈群”一词的由来。]

中心扩张的圈代数有远为丰富的结构与理论。如上所见，它们由第二上同调群 $H^2(L\mathfrak g,\C)$ 分类；对单李代数 $\mathfrak g$，已知该群是一维的。相应的上循环由公式 (1.3-3) 给出。它们依赖于 $\mathfrak g$ 上记作 $\kappa$ 的不变内积。我们把中心扩张记作 $\gghat^{\mathrm{pol}}_\kappa$。\transnote{公式 (1.3-3) 的具体形式是 $c(x\otimes f,y\otimes g)=\kappa(x,y)\operatorname{Res}_{t=0}(f\,\mathrm dg)$，其中 $x,y\in\mathfrak g$、$f,g\in\C((t))$，$\kappa$ 是 $\mathfrak g$ 上 $\mathrm{ad}$-不变的双线性型（对单 $\mathfrak g$ 唯一到常数倍）。不变性保证 Jacobi 型条件，而 $\operatorname{Res}(f\,\mathrm dg)=-\operatorname{Res}(g\,\mathrm df)$ 给出反称性；$H^2$ 的一维性意味着所有一维中心扩张都是某个 $\kappa$ 的倍数。[来源：原书 p.385（引式 (1.3-3)，见原书 §1.3）；Kac 1990, Ch.7（仿射代数作为圈代数的中心扩张）。]}

上述构造还有一个形式版本：把洛朗多项式代数 $\C[t,t^{-1}]$ 换成形式洛朗级数代数 $\C((t))$。这给出形式圈代数 $\mathfrak g\otimes\C((t))=\mathfrak g((t))$。它的中心扩张由其第二上同调群参数化；在适当的连续性条件下，对单李代数 $\mathfrak g$ 这也是一维向量空间。中心扩张同样由指定 $\mathfrak g$ 上的不变内积 $\kappa$ 给出。李括号由与之前相同的公式给出——注意留数仍然良定义。我们把相应的李代数记作 $\gghat_\kappa$。这就是与 $\mathfrak g$（及 $\kappa$）相伴的仿射 Kac--Moody 代数。\transnote{附录没有给出 $\kappa$ 的归一化，但第 2、3 章实际取 $\kappa_0$ 使长根的平方长度为 $2$（即 $(\theta|\theta)=2$，$\theta$ 为最高根），任一不变型写作 $\kappa=k\kappa_0$；对偶 Coxeter 数 $h^\vee=1+\langle\rho,\theta^\vee\rangle$（$\mathfrak{sl}_2$ 为 $2$，$\mathfrak{sl}_n$ 为 $n$）。归一化之所以重要：临界水平 $\kappa_c=-h^\vee\kappa_0$ 即 $k=-h^\vee$，而 Virasoro 中心荷 $c=k\dim\mathfrak g/(k+h^\vee)$ 的分母 $k+h^\vee$ 只在上述归一化下成立；换归一化会使 $k$ 与 $h^\vee$ 同步改变，公式须整体换算。[来源：原书 p.38 式 (2.1-10)--(2.1-11)、p.63--p.64；Kac 1990, Ch.6（标准化不变型、对偶 Coxeter 数与 $\rho$）。]}

% [Tutor Note: 原书笔误（p.385）。原文 “This gives is the formal loop algebra” 多写了一个 “is”；译文按语义只译一次，此处登记。]

% [Tutor Note: 背景（p.385）。$\gghat^{\mathrm{pol}}_\kappa$ 与 $\gghat_\kappa$ 的关系是：前者基于多项式圈，后者基于形式圈；两者中心扩张都由同一不变内积 $\kappa$ 给出，但后者才有“留数”与连续对偶等分析结构，因而成为本书主角。$\kappa$ 与 $\kappa_c$ 的临界水平、对偶 Coxeter 数 $h^\vee$、以及中心 $Z(\gghat)$ 同构于 opers 上函数（Feigin--Frenkel）都在这条线上展开。]
